
Der Kompetenz-Fit beschreibt die Passung zwischen den Kompetenzanforderungen eines ERP-Systems und den individuellen sowie 
organisatorischen Fähigkeiten eines KMU. Während der Ressourcen-Fit vor allem die Verfügbarkeit von Budget, Zeit und Personal 
betrachtet, bezieht sich der Kompetenz-Fit auf die Frage, ob das Unternehmen über ausreichendes Wissen verfügt, um ein ERP-System 
auszuwählen, einzuführen, zu nutzen und weiterzuentwickeln. Ein hoher Kompetenz-Fit liegt somit vor, wenn die beteiligten Personen 
die fachlichen, technischen und organisatorischen Anforderungen eines ERP-Projekts verstehen und bewältigen können.

ERP-Projekte erfordern Kompetenzen, die über reines IT-Wissen hinausgehen. Dazu gehören insbesondere ein Verständnis der eigenen 
Geschäftsprozesse, Kenntnisse über Datenstrukturen, die Fähigkeit zur Formulierung von Systemanforderungen sowie 
Projektmanagement- und Veränderungskompetenzen. Mitarbeitende aus den Fachbereichen müssen beispielsweise beurteilen können, 
welche Anforderungen tatsächlich notwendig sind und wie bestehende Arbeitsabläufe künftig im ERP-System unterstützt werden sollen. 
Gleichzeitig müssen technische Verantwortliche die Systemarchitektur, Schnittstellen, Datenmigration und mögliche Anpassungen 
einschätzen können \parencite{Alaskari2021,Leyh2015}.

Gerade in KMU ist diese Kompetenzbasis häufig nur begrenzt vorhanden. Viele Unternehmen verfügen über kleine IT-Abteilungen oder 
keine spezialisierten Mitarbeitenden für ERP- und Digitalisierungsprojekte \parencite{Faruque2024}. Dadurch entsteht häufig eine starke 
Abhängigkeit von externen Beratungen, Implementierungspartnern oder Softwareanbietern. Externe Unterstützung ist grundsätzlich nicht problematisch, 
kann jedoch kritisch werden, wenn das Unternehmen seine eigenen Anforderungen nicht präzise beschreiben oder externe Empfehlungen 
nicht angemessen bewerten kann. In diesem Fall besteht die Gefahr, dass Entscheidungen überwiegend durch externe Akteure getroffen 
werden und das ERP-System nicht ausreichend auf die tatsächlichen Bedürfnisse des Unternehmens abgestimmt wird \parencite{Leyh2015,Leyh2017}.

%Ein hoher Kompetenz-Fit setzt daher nicht voraus, dass ein KMU sämtliche technischen Kenntnisse intern aufbauen muss. Entscheidend ist 
%vielmehr, dass das Unternehmen eine kompetente Auftraggeberrolle einnehmen kann. Dazu gehört, dass zentrale Ansprechpartnerinnen und 
%Ansprechpartner vorhanden sind, die die Fachbereiche vertreten, Anforderungen bündeln, Entscheidungen vorbereiten und die 
%Zusammenarbeit mit externen Partnern steuern können. Besonders relevant sind sogenannte Key User, die sowohl die Anforderungen 
%ihres Fachbereichs kennen als auch ausreichend Verständnis für die Funktionsweise des ERP-Systems entwickeln. Sie können als 
%Vermittler zwischen Fachabteilungen, IT, Management und externen Dienstleistern fungieren.

Neben der Kompetenz zur Planung und Einführung ist auch die Fähigkeit zur späteren Nutzung des ERP-Systems entscheidend. Mitarbeitende 
müssen verstehen, welche Auswirkungen ihre Eingaben und Arbeitsschritte auf andere Unternehmensbereiche haben. Werden Daten 
beispielsweise fehlerhaft gepflegt oder Prozessschritte nicht korrekt durchgeführt, können sich die Folgen über mehrere Bereiche 
hinweg auswirken. ERP-Kompetenz umfasst daher auch ein Verständnis für integrierte Daten- und Prozesszusammenhänge. Dies 
unterscheidet die Arbeit mit einem ERP-System von der Nutzung isolierter Einzelanwendungen \parencite{Leyh2015,Alaskari2021}.

Eine zentrale Maßnahme zur Entwicklung dieser Kompetenzen sind Schulungen. Schulungen sollten nicht ausschließlich die technische 
Bedienung einzelner Funktionen vermitteln, sondern sich an den konkreten Rollen und Aufgaben der Mitarbeitenden orientieren 
\parencite{Noudoostbeni2009,Faruque2024}. 

Eng mit Kompetenzen verbunden ist die Akzeptanz der Mitarbeitenden. Ein ERP-System kann technisch korrekt implementiert sein und 
dennoch nur begrenzt wirksam werden, wenn es im Arbeitsalltag nicht angenommen wird. Fehlende Akzeptanz zeigt sich beispielsweise darin, dass 
Mitarbeitende weiterhin eigene Excel-Dateien nutzen, Daten außerhalb des Systems pflegen oder neue Prozessschritte umgehen 
\parencite{Leyh2015,Leyh2019}. Dadurch entstehen erneut parallele Informationsbestände, die die angestrebte Integration und Transparenz 
beeinträchtigen können.

%Akzeptanz entsteht nicht allein durch Schulung, sondern auch durch frühzeitige Einbindung und nachvollziehbare Kommunikation. 
%Mitarbeitende sollten verstehen, warum sich Prozesse verändern, welche Vorteile das System für ihre tägliche Arbeit bringen kann 
%und welche Erwartungen mit der Einführung verbunden sind. Werden Fachbereiche frühzeitig an Prozessanalysen, Tests oder 
%Schulungskonzepten beteiligt, kann dies dazu beitragen, Widerstände zu reduzieren und praktisches Wissen in die Systemgestaltung 
%einzubringen. Kompetenz-Fit umfasst somit auch die Fähigkeit eines Unternehmens, Veränderungsprozesse aktiv zu begleiten.

Darüber hinaus ist der langfristige Wissenserhalt relevant. Wenn zentrale Kenntnisse ausschließlich bei externen Beratern oder 
einzelnen Mitarbeitenden liegen, entsteht nach der Einführung ein erhebliches Abhängigkeitsrisiko. Für einen nachhaltigen Kompetenz-Fit 
sollte Wissen daher schrittweise im Unternehmen verankert werden. Dies kann durch dokumentierte Entscheidungen, interne Ansprechpersonen, 
Key-User-Konzepte und einen gezielten Wissenstransfer während der Projektlaufzeit unterstützt werden. Auf diese Weise kann ein KMU das 
ERP-System auch nach dem Go-Live weiterentwickeln und an neue Anforderungen anpassen \parencite{Alaskari2021,Jha2008,Leyh2015}.

Für die Bewertung des Kompetenz-Fits sind insbesondere die vorhandene ERP- und IT-Erfahrung, das Prozessverständnis der Fachbereiche, 
die Fähigkeit zur Steuerung externer Partner, die Qualität des Schulungs- und Change-Konzepts sowie die Akzeptanz der späteren 
Anwenderinnen und Anwender relevant. Ein niedriger Kompetenz-Fit liegt vor, wenn Anforderungen nur unklar formuliert werden können, 
Wissen überwiegend extern liegt, Schulungen fehlen oder Mitarbeitende die neuen Prozesse nicht akzeptieren. Ein hoher Kompetenz-Fit 
liegt dagegen vor, wenn fachliche und technische Kompetenzen sinnvoll zusammenwirken, Schlüsselpersonen vorhanden sind, Wissen im 
Unternehmen aufgebaut wird und die Nutzung des ERP-Systems aktiv unterstützt wird.

%Zusammenfassend beschreibt der Kompetenz-Fit die Fähigkeit eines KMU, ein ERP-System nicht nur technisch einzuführen, sondern es auch 
%organisatorisch zu verstehen, im Arbeitsalltag zu nutzen und langfristig weiterzuentwickeln. Damit bildet er eine wesentliche 
%Voraussetzung dafür, dass die mit ERP-Systemen verbundenen Integrations- und Transformationspotenziale tatsächlich realisiert 
%werden können.
